\subsection{自反 Fréchet 空间的对偶}

命题 4. --- 设 E 是一个自反 Fréchet 空间。E 的强对偶 \(E_{b}^{\prime}\) 是 Banach 空间的一个序列的归纳极限。

设 \((\mathrm{V}_{n})_{n\in\mathbf{N}}\) 是 E 中凸、均衡且闭的零邻域的一个递减序列，使得 E 的每个零邻域都包含某个 \(V_{n}\)。设 \(A_{n}\) 是 \(V_{n}\) 在 \(E'\) 中的极集。则 \(A_{n}\) 是凸、均衡且对于 \(\sigma(\mathrm{E}',\mathrm{E})\) 紧的；由 III 第 8 页的推论，空间 \(E_{A_{n}}'\) 是 Banach 空间。我们将证明 \(E_{b}'\) 是诸空间 \(E_{A_{n}}'\) 的归纳极限；换言之，\(E'\) 的每个吸收诸 \(A_{n}\) 中每一个的凸、均衡子集 U 都是 \(E_{b}'\) 中 0 的一个邻域。对每个 \(n\in\mathbf{N}\)，选取实数 \(\lambda_{n}>0\) 使得 \(\lambda_{n}A_{n}\subset U\)。设 \(B_{n}\) 是集 \(\bigcup_{0\leqslant i\leqslant n}\lambda_{i}A_{i}\) 的凸包；令 \(V=\bigcup_{n}B_{n}\)，则 \(V\subset U\)。对每个 \(n\in\mathbf{N}\)，集 \(B_{n}\) 是凸、均衡且对于 \(\sigma(\mathrm{E}',\mathrm{E})\) 紧的（II，第 14 页，命题 15）。

现在我们将证明 \(\frac{1}{2}\mathrm{V}^{\circ \circ}\subset \mathrm{V}\)。设 \(x\in \mathrm{E}_b^\prime -\mathrm{V}\)；对每个 \(n\in \mathbf{N}\)，我们有 \(x\notin \mathbf{B}_n\)，且由于 \(\mathbf{B}_n\) 对于 \(\sigma (\mathrm{E}',\mathrm{E})\) 闭，存在元素 \(y_{n}\)（在 \(\mathbf{B}_n^\circ\) 中的）使得 \(\langle y_n, x \rangle = 1\)（II，第 38 页，命题 4）。由于 E 是自反的，E 的每个有界子集对于 \(\sigma(E, E')\) 相对紧（IV，第 16 页，定理 2）。由 \(B_n\) 的定义，我们有

\[
\lambda_ {i} y _ {n} \in \mathrm{V} _ {i} \quad \text { for   all } \quad n \geqslant i  ,\tag{3}
\]

从而序列 \((y_{n})\) 有界。设 \(y\) 是 \((y_{n})\) 对于拓扑 \(\sigma(\mathrm{E},\mathrm{E}')\) 的一个极限点。我们有 \(y\in\mathrm{V}^{\circ}=\bigcap_{n}\mathrm{B}_{n}^{\circ}\) 且 \(\langle y,x\rangle=1\)。于是 \(x\notin\frac{1}{2}\mathrm{V}^{\circ\circ}\)，从而有包含关系 \(\frac{1}{2}\mathrm{V}^{\circ\circ}\subset\mathrm{V}\)，更加有 \(\frac{1}{2}\mathrm{V}^{\circ\circ}\subset\mathrm{U}\)。

由于 \(E_{b}^{\prime}\) 的每个有界子集都包含在某个集 \(A_{n}\) 中，集 \(V = \bigcup_{n} B_{n}\)

吸收 \(E_{b}^{\prime}\) 的每个有界子集。于是，\(V^{\circ}\) 在 E 中有界，从而 \(\frac{1}{2}V^{\circ\circ}\) 是 \(E_{b}^{\prime}\) 中 0 的一个邻域。更加地，U 是 \(E_{b}^{\prime}\) 中 0 的一个邻域。

推论. --- 自反 Fréchet 空间的强对偶是有界型的且是桶型的。

Banach 空间的归纳极限按定义是有界型的。此外，Banach 空间是桶型的（III，第 25 页，推论），且桶型空间的每个归纳极限都是桶型空间（III，第 25 页，推论 3）。

